\subsection{A periodic triangular source opposite a right-zero corner}

\begin{lemma}[Proved: a triangular loop algebra on a two-term complex]
\label{CP2:A-triangular-loop-two-term-amplitude}
Let $A$ be the four-dimensional rational algebra with basis
$1,h,x,z$, with identity $1$, and with products
\[
 h^2=h,\qquad hz=z,
\]
all other products of two elements of $\{h,x,z\}$ being zero.
Let $k$ be the one-dimensional left module on which $1$ acts as
the identity and $h,x,z$ act as zero.  If $D$ is a complex of
finite-dimensional right $A$ modules in degrees zero and one,
and $D\otimes_A^{\mathbf L}k$ has bounded homology, then
\[
 H_j(D\otimes_A^{\mathbf L}k)=0\qquad(j\ge2).
\]
\end{lemma}
\begin{proof}
Put $p=1-h$.  The projective left modules $Ap$ and $Ah$ have
bases $p,x,z$ and $h$, respectively.  In particular $x,z\in Ap$,
$px=x$, $hz=z$, $xh=zh=0$ and $pz=0$.
The augmentation $Ap\to k$ takes $p$ to $1$ and kills $x,z$.
Its kernel is the span of $x,z$.

The following is an exact projective resolution of $k$:
\[
 P_0=Ap,\qquad P_j=Ap\oplus Ah\quad(j\ge1),
\]
\[
 d_1(u,v)=ux+vz,\qquad
 d_j(u,v)=(ux+vz,0)\quad(j\ge2).
\]
Right multiplication by $x$ maps $Ap$ onto $\mathbb Qx$ and
has kernel $\mathbb Q\{x,z\}$.  Right multiplication by $z$
maps $Ah$ isomorphically onto $\mathbb Qz$.  Thus $d_1$ has
image $\mathbb Q\{x,z\}$ and kernel
$\mathbb Q\{(x,0),(z,0)\}$.  Every later differential has
the same kernel and has that kernel as its image in the preceding
term.  These statements prove exactness in every degree.

Write $\partial:D_1\to D_0$.  For $i=0,1$, put
\[
 E_i=D_i p\oplus D_i h,\qquad
 \psi_i:E_i\longrightarrow D_i p,\qquad
 \psi_i(u,v)=ux+vz,
\]
and let $\phi_i:E_i\to E_i$ be $\phi_i(u,v)=(\psi_i(u,v),0)$.
The target of $\psi_i$ is correct because $xp=x$ and $zp=z$.
Write $\partial_E:E_1\to E_0$ for the direct sum of the two
restrictions of $\partial$; they are defined because $\partial$
is $A$ linear.

The total tensor complex $C=D\otimes_A P$ has
\[
 C_0=D_0p,\qquad C_1=E_0\oplus D_1p,\qquad
 C_j=E_0\oplus E_1\quad(j\ge2).
\]
With the usual sign $(-1)^i$ on the resolution differential
in the row $D_i$, all differentials $d_j$ for $j\ge3$ are the
same map
\[
 T:E_0\oplus E_1\longrightarrow E_0\oplus E_1,
 \qquad T(u,v)=(\phi_0u+\partial_Ev,-\phi_1v).
\]
At the degree-two transition one instead has
\[
 d_2(u,v)=(\phi_0u+\partial_Ev,-\psi_1v)
       \in E_0\oplus D_1p.
\]
Let $\iota:C_1\hookrightarrow E_0\oplus E_1$ insert zero in
the $D_1h$ coordinate.  This is an injection and satisfies
$\iota d_2=T$.  The omitted coordinate is therefore identically
zero, rather than an additional constraint.  Consequently
\[
 \ker d_2=\ker T,\qquad \operatorname{im}d_3=\operatorname{im}T,
\]
and, for every $j\ge2$,
\[
 H_j(C)=\ker T/\operatorname{im}T.
\]
The displayed projective resolution computes the derived tensor
for the finite complex $D$.  If its homology is bounded, this
repeated vector space must be zero.  This proves the lemma,
including the degree-two endpoint.
\end{proof}

\begin{theorem}[Proved: the first surviving eight-element source opposite $R_3$]
\label{CP2:A-eight-source-five-right-zero}
Let $K=\{e_{ij}:i,j\in\{0,1\}\}$ be the rectangular band,
$e=e_{00}$, and let $M=K\amalg\{1,a,b,c\}$ have identity $1$,
band products $e_{ij}e_{kl}=e_{il}$, and
\[
 ak=bk=ck=ek,\qquad ka=kb=kc=ke\quad(k\in K).
\]
Among the products of two elements of $\{a,b,c\}$, let
$c^2=c$ and $cb=b$, and let every other such product be $e$.
Let $N=R_3$ be the identity-adjoined two-element right-zero band.
For every finite nonempty left-$N$, right-$M$ biset $S$, bounded
rational homology of the collage implies that all its rational
degree-two products are zero.  There is no bound on $|S|$ or on
its second Betti number.

This is source $5$ in the frozen positive order-eight catalog,
so its entire contribution to the eighteen-arrow budget $(8,3,7)$
is excluded.
\end{theorem}
\begin{proof}
We use the natural graph localization and verify the additional
algebraic hypothesis explicitly.  The left $R_3$ action has the
equivariant graph model of
Corollary~\ref{CP2:A-baseline-right-zero-three-six}: the vertices
are the common image of the two band retractions, and each state
outside that image supplies an edge between its two images.
Every commuting right $M$ actor sends an edge, with its parameter,
to an edge or a vertex.  This is the natural augmented comparison
used in that corollary; taking the right-action bar preserves it.
Since $BR_3$ is contractible, the collage is rationally equivalent,
by an actual comparison respecting products, to the graph cofiber
\[
 X=\operatorname{hocofib}(\Gamma_{hM}\longrightarrow BM).
\]

Put $R=\mathbb QM$, $A=R/\mathbb QK$, and let $V$ be the
left row-augmentation module.  We recall the graph localization
calculation from the proof of
Theorem~\ref{CP2:A-projective-local-row-graph-obstruction}; its
calculation before the final algebraic case split has no localness
or projectivity assumption on $V$.
If $\Gamma e_{00}$ and $\Gamma e_{01}$ have a common vertex,
that vertex is fixed by every band element and hence every element
of $M$, since $em\in K$.  The resulting bridge section makes
every second cohomology class of the cofiber relative to $BM$.
The actual suspension quotient then kills all degree-two products.

Otherwise the two image graphs have disjoint vertex and edge
sets.  Set
\[
 C=[C_1(\Gamma;\mathbb Q)\longrightarrow
                  \ker(C_0(\Gamma;\mathbb Q)\to\mathbb Q)]
\]
in degrees one and zero.  Multiplication embeds
$Ce\otimes_{\mathbb Q}eR$ in $C$ with image $C\mathbb QK$:
its two column copies have disjoint supports, and in degree zero
both have augmentation zero.  Hence
\[
 D=C/(Ce\otimes_{\mathbb Q}eR)
\]
is a two-term right $A$ complex.  The removed subcomplex consists
of projective right $R$ modules, so its tensor contributes only
in degrees zero and one.  The row sequence
$0\to V\to Re\to\mathbb Q\to0$ and $De=0$ supply a shift.
The exact sequence $0\to\mathbb QK\to R\to A\to0$, with
$\mathbb QK=Re\otimes eR$, identifies the resulting derived
tensor over $R$ with that over $A$.  Together with the cofiber
shift, this gives the natural localization identity
\[
 H_j(X;\mathbb Q)=H_{j-2}(D\otimes_A^{\mathbf L}V)
                    \qquad(j\ge4).
\]
In particular, boundedness of $X$ implies boundedness of the
derived tensor on the right; its finitely many lower degrees
do not affect that conclusion.

In the literal contracted algebra, take
\[
 h=\overline c,\qquad x=\overline a,\qquad z=\overline b.
\]
Its table is exactly the algebra of the preceding lemma:
$h^2=h$, $hz=z$, and all its other products between
$h,x,z$ are zero.  Each of $a,b,c$ acts constantly on the two
minimal rows, so $V$ is exactly the character module $k$ of
that lemma.  The lemma forces
$H_2(D\otimes_A^{\mathbf L}V)=0$, and the localization identity
therefore gives $H_4(X;\mathbb Q)=0$.  This also kills every
degree-two product, proving the claim in the second case.

The separate literal binding uses catalog identity $0$, band
$\{1,2,3,4\}$ with $e=1$, and $(a,b,c)=(5,6,7)$.  Every
one of its sixty-four products and the contracted algebra table
are reconstructed in
\path{acyclic/outputs/eight_source_five_literal_binding.json}.
Thus the stated finite source identification uses no unverified
classification of the connecting actions.
\end{proof}
